%!TEX root = ../dissertation.tex
\chapter{Conclusion}
In this thesis, we presented an in-depth proof of Vizing's theorem on the chromatic index of a simple graph. This compilation of proof sketches, published algorithms, and some novel lemmas creates a single source for all the required components in a formal proof of Vizing's theorem. We implemented the contents of this paper in the proof assistant Rocq. Our main contribution is the definition of an edge coloring and related basic operations within the framework of Rocq's graph theory library. Once we became familiar with typical Rocq and SSReflect practices as well as the nuances of the graph theory library, the challenge became implementing edge colorings in a way that felt natural and would be the simplest to interact with. With a definition of edge colorings centered around mappings, we formalized and verified properties about fans, rotation, Kempe chains, and Kempe swaps. Maximal length fans and Kempe chains as alternating paths demanded we familiarize ourselves with Fixpoints and recursion to ensure termination, a strict requirement of Rocq functions. Due to anticipated complications creating maximal Kempe chains, we researched alternative representations and discovered we could sidestep infinite loops entirely with a subgraph representation. We then needed to prove more general lemmas relating connected paths and vertex degree before applying the chains to our final proof. We concluded the project with a verified proof of Vizing's theorem for both Kempe chains as alternating paths and Kempe chains as subgraph components. 

\subsection{Future Work}
While we completed our original goal to implement and verify Vizing's theorem in the proof assistant Rocq, there are many opportunities to extend the project. Firstly, a few admitted lemmas are still scattered within the code as discussed in Section~\ref{section admit it}, so the next immediate goal is to complete the corresponding proofs. Secondly, we believe this project would be a valuable contribution to the graph theory library created by Doczkal and Pous. The next step here is to connect with the library's maintainers. After this, we imagine we will need to edit our library to ensure it matches stylistically, and some lemmas may be made more or less general. We ideally could work in stages, first integrating edge colorings and then progressing towards Kempe chains and fans before eventually contributing to Vizing's theorem itself. Finally, on a related note, there are several formalized lemmas that could be generalized or be starting points for other graph theory results. We have shown that we can extend an edge coloring when a new edge is added to a graph and preserve the properness of the edge coloring. We could support adding any set of non-adjacent edges, in contrast to only a single edge, and still preserve the properness of an edge coloring. Additionally, some aspects of this library, including edge sets at a vertex as discussed in Section~\ref{section d&d}, could apply to both directed and simple graphs, but we currently only reason over simple graphs. Vizing's theorem itself can be generalized to multigraphs, and proving this would be a worthwhile endeavor. Beyond this, there is a plethora of literature surrounding edge colorings and chromatic index. We hope that future work will continue to improve upon our implementation and one day any edge coloring result will be able to be formally stated in a proof assistant.

\label{conclusion}